\subsection{Reduced Ground Instances, Basic Proofs, and the Hereditary Restrictions}
\label{sec:p12-eu-hereditary-source}

The source's Definitions 4.3--4.6 and Lemmas 4.7--4.9 supply a
specific reduced-instance route, distinct from the later arbitrary-theory
completion construction. Its objects and proof-length measure are retained
here, with explicit corrections to the individual-step descent argument.
\cite[pp.~489--495]{BaaderSnyder2001}

\begin{hbdefinition}{def:p12-eu-h-ground-instances}{The source's reduced substitutions and ground-instance sets}
Fix nontrivial E and a reduction ordering \(\succ\), total on the
chosen ground signature. An R-reduced substitution has every image
\(x\theta\), for \(x\in\operatorname{Dom}(\theta)\), R-normal.
Put
\[
\begin{aligned}
 \operatorname{Gr}(E)&=\{l\rho\approx r\rho:
     l\approx r\in E\cup E^{-1},\ l\rho,r\rho\text{ ground}\},\\
 R&=\operatorname{Gr}^{\succ}(E)
   =\{l\rho\to r\rho:l\rho\approx r\rho\in\operatorname{Gr}(E),
                                               l\rho\succ r\rho\},\\
 R_E&=\{l\rho\to r\rho\in R:
          \rho\text{ is R-reduced on the identity variables}\}.
\end{aligned}
Use the generating substitution only on
\(\operatorname{Var}(l,r)\); arbitrary unrelated assignments are
not part of the instance's provenance. The source's ``entire set''
for reducedness is R. Thus \(R_E\) is a defined subset of R,
not an assumption that every ground instance is reduced.
\cite[Definitions 4.3--4.4, pp.~489--490]{BaaderSnyder2001}
\end{hbdefinition}

\begin{hblemma}{thm:p12-eu-h-minimal-instance}{proved=annex:p12.n-c-035.5.h-minimal-instance}{Minimal applicable pairs really have reduced generating images}
In a finite term admitting an R-redex, choose an applicable pair
\((l\rho,r\rho)\) lexicographically minimal under \(\succ\).
Every generating image is R-normal, including variables occurring
only on the identity's right side. Consequently that redex is an
\(R_E\)-redex. R and \(R_E\) have precisely the same normal
terms. The choice is an existence construction, not an effective
enumeration or a claimed finite branching bound.
\end{hblemma}

\begin{hbdefinition}{def:p12-eu-h-basic-proof}{Based steps and the evolving source skeleton}
A step from \(q\theta\) is based on q when its redex is at a
nonvariable position in q. With a fresh rule presentation
\(l\rho\to r\rho\), replace the corresponding skeleton
occurrence by r, retaining its variables uninstantiated. The ground
successor is \(q[r]\theta\rho\), and the next step must be based
on this new skeleton. A finite based sequence is reflexive or repeats
these operations. A basic rewrite proof consists of based sequences
from \(s\theta\) and \(t\theta\) to their common normal form.
Normality of the final common term is part of ``rewrite proof'';
a zero-step equality at a reducible term is not such a proof.
\cite[Definition 4.6 and preceding convention, pp.~488,490]{BaaderSnyder2001}
\end{hbdefinition}

\begin{hbtheorem}{thm:p12-eu-h-basic-normalization}{proved=annex:p12.n-c-035.5.h-basic-normalization}{The original reduced-instance normalization and basic-proof equivalence}
Suppose R is ground convergent and \(s\theta\) is ground, with
\(\theta\) \(R_E\)-reduced. Every R-normalization of it has the
same endpoint as a finite based normalization using only \(R_E\).
For ground \(s\theta,t\theta\) with R-reduced \(\theta\),
\(s\theta=_E t\theta\) holds exactly when a basic rewrite proof
to a common normal form exists. It may be chosen using \(R_E\).
\cite[Lemma 4.7 and Corollary 4.8, p.~491]{BaaderSnyder2001}
\end{hbtheorem}

\begin{hblemma}{thm:p12-eu-h-basic-transport}{proved=annex:p12.n-c-035.5.h-basic-transport}{Literal solving does not destroy a reduced basic witness}
Let a finite \(R_E\)-basic proof be witnessed by reduced ground
\(\alpha\). If a finite syntactic substitution \(\mu\) satisfies
\(\mu\alpha=\alpha\) on its relevant variables, applying \(\mu\)
to the proof's initial skeletons preserves a basic proof with exactly
the same ground steps. Future identity variants are chosen fresh for
\(\mu\). In particular the minimum basic-proof length cannot
increase under a finite literal matching solve preserving that witness.
\end{hblemma}

\begin{hbdefinition}{def:p12-eu-h-source-measure}{The printed four-tuple and its explicit reduced-proof reading}
For \(P;S\) and a reduced witness \(\theta\), the source uses
\((m,n_1,n_2,n_3)\): m totals the lengths of minimal basic
normal-form rewrite proofs of the equations in \(P\theta\);
\(n_1\) counts distinct variables in equations whose instantiated
sides are literally the same normal term; \(n_2\) counts their
term-symbol occurrences; \(n_3\) counts their nonvariable-left,
variable-right orientations. Counts include equation multiplicity,
and may be zero; the order is lexicographic on natural numbers.

For the corrected reduced-instance proof, write \(m_{R_E}\) for
that minimum with all steps explicitly restricted to \(R_E\).
It exists by \cref{thm:p12-eu-h-basic-normalization}. The remaining
three coordinates are the source's same counts. The distinction
matters: a shortest unrestricted R-proof may erase a nonnormal image.
The induction step below is LP followed by its mandatory first
Decomposition and a full finite syntactic solve of the introduced
matching equation, not a raw LP step alone.
\cite[Lemma 4.9 proof, pp.~491--492]{BaaderSnyder2001}
\end{hbdefinition}

\begin{hbexample}{def:p12-eu-h-measure-failures}{Three different claims about shortest proofs}
For \(E=\{f(a)\approx a\}\), input \(f(a)=^{?}a\) has
minimum normal-form proof length one. Raw LP creates
\(a=^{?}a\) and \(f(a)=^{?}f(a)\); the latter needs two
steps to its common normal form, so m increases to two.

For \(E=\{f(x)\approx x,f(f(x))\approx x\}\), input
\(f(f(a))=^{?}a\) has a one-step reduced-instance proof. LP
with the second identity followed only by its first Decomposition
leaves \(f(x)=^{?}f(a)\), witnessed by \(x=a\); its two
normalization legs still total two. Full literal matching solving
is needed before asserting the source tuple's strict decrease.

Finally take \(E=\{f(x)\approx a,g(b)\approx c\}\) and
\(f(g(b))=^{?}a\). Its one-step unrestricted basic proof uses
\(x\mapsto g(b)\), which is not reduced. A reduced-instance
proof instead reduces \(g(b)\) to c and then \(f(c)\) to a.
This is why existence of a basic proof does not make its unrestricted
shortest version a reduced-instance proof.
\end{hbexample}

\begin{hbtheorem}{thm:p12-eu-h-minimum-blocks}{proved=annex:p12.n-c-035.5.h-minimum-blocks}{Lemma 4.9 with the source proof-length objects and a justified finite block}
Let nontrivial E have ground-convergent R in the declared parameter
signature, and let reduced \(\theta\) E-solve
\(P;\varnothing\) with \(P\theta\) ground. For open images,
first freeze their finite parameter set and require these hypotheses
in that fresh-constant extension, as in
\cref{def:p12-eu-narrowing-standing}. There is a finite
G trace to \(\varnothing;S\), with
\(x\theta=x\sigma_S\lambda\) on
\(\mathcal X=\operatorname{Var}(P)\).
Each selected normal-equation U step decreases
\((m_{R_E},n_1,n_2,n_3)\), and each selected reduced-instance
LP/full-matching block strictly decreases its first coordinate.
Fresh variables introduced by LP do not change old witness images.
This is a proof using the source's minimal basic-proof route with
the repairs stated above, not the earlier ground-term-multiset route.
\cite[Lemma 4.9, pp.~491--492]{BaaderSnyder2001}
\end{hbtheorem}

\begin{hbdefinition}{def:p12-eu-h-hereditary-tests}{The two source hereditary constraints}
Full nonvariable overlap compatibility requires that, at every position
where both l and u have a nonvariable symbol, their symbols and arities
agree. It does not require variables to have matching names, nor does
it replace the literal matching problem by E-equality.
The historical basic constraint records occurrence provenance:
constructor occurrences introduced by Variable Elimination inside
a substituted image are ineligible for subsequent LP, even though
they are physically present in the raw instantiated problem. Original
and freshly inserted identity-template constructor occurrences remain
eligible. Decomposition inherits the occurrence marks; syntactic
solving is not itself prohibited at marked occurrences.
These are the two restrictions described in the source's \S4.2.
\cite[p.~495]{BaaderSnyder2001}
\end{hbdefinition}

\begin{hbtheorem}{thm:p12-eu-h-convergent-hereditary}{proved=annex:p12.n-c-035.5.h-convergent-hereditary}{Both hereditary restrictions preserve the reduced-instance trace}
Under the exact ground-convergence/reduced-witness hypotheses of
\cref{thm:p12-eu-h-minimum-blocks}, its finite trace can satisfy
both hereditary constraints. Each LP match has a literal ground
witness, proving compatibility at every shared nonvariable position.
Every substituted copy retained in P has a normal ground value;
therefore the chosen next rewrite cannot be inside it. This establishes
actual occurrence eligibility, not a definition naming an unspecified
``basic lifting lemma.''
\end{hbtheorem}

\begin{proofboundary}{def:p12-eu-h-arbitrary-refinement-boundary}{The arbitrary-theory refinement is a separate remaining source argument}
The source's \S4.2 refers the arbitrary-theory refinements to
Dougherty--Johann 1992, Socher-Ambrosius 1994, and Moser 1993, and
explicitly does not develop their detailed proofs. The theorem just
proved is the source reduced-instance/ground-convergent regime; it
does not erase the arbitrary-theory refinement obligation.
The existing arbitrary-theory G theorem alone cannot discharge it:
critical-pair simulation may introduce function occurrences through
its unifier, whose historical eligibility differs from a primitive
critical-identity template. A proof must actually transport those
marks and the full-overlap tests through every nested simulation,
or reconstruct the identified refinement's own proof method.
\cite[\S4.2, p.~495]{BaaderSnyder2001}
\end{proofboundary}
